\documentclass[12pt]{article}
\usepackage[margin=1in]{geometry}
\usepackage[utf8]{inputenc}
\usepackage{latexsym, amsfonts, enumitem, graphicx, environ,enumitem, fancyhdr, color, amssymb, amsthm, amsmath, mathtools, tikz, nicematrix}
\usetikzlibrary{patterns}
\usetikzlibrary{decorations.pathreplacing}
\pagestyle{fancy}
\setlength{\headheight}{65pt}
\newenvironment{problem}[2][Problem]{\begin{trivlist}
    \item[\hskip \labelsep {\bfseries #1}\hskip \labelsep {\bfseries #2.}]}{\end{trivlist}}
\DeclareMathOperator{\Ima}{Im}

\usepackage{hyperref}
\usepackage[capitalize]{cleveref}

\newtheorem{thm}{Theorem}[section]
\newtheorem{lem}[thm]{Lemma}
\newtheorem{cl}[thm]{Claim}
\newtheorem{prop}[thm]{Proposition}
\newtheorem{cor}[thm]{Corollary}
\newtheorem{conj}[thm]{Conjecture}
\newtheorem{obstruction}[thm]{Obstruction}

\usepackage{mdframed}

\theoremstyle{definition}
\newtheorem{definition}[thm]{Definition}
\newtheorem{remark}[thm]{Remark}
\newtheorem{question}{Question}
\newtheorem{example}[thm]{Example}
\newtheorem{exercise}[thm]{Exercise}
\newtheorem{properties}[thm]{Properties}

\usepackage{tcolorbox}
\tcbuselibrary{theorems,skins,breakable}

\tcolorboxenvironment{example}{
enhanced jigsaw,colframe=cyan,interior hidden,
breakable,%before skip=10pt,after skip=10pt
}

\tcolorboxenvironment{thm}{
enhanced jigsaw,colframe=red,interior hidden,
breakable,%before skip=10pt,after skip=10pt
}

\tcolorboxenvironment{prop}{
enhanced jigsaw,colframe=red,interior hidden,
breakable,%before skip=10pt,after skip=10pt
}

\tcolorboxenvironment{properties}{
enhanced jigsaw,colframe=red,interior hidden,
breakable,%before skip=10pt,after skip=10pt
}

\tcolorboxenvironment{lem}{
enhanced jigsaw,colframe=red,interior hidden,
breakable,%before skip=10pt,after skip=10pt
}
\tcolorboxenvironment{cor}{
enhanced jigsaw,colframe=red,interior hidden,
breakable,%before skip=10pt,after skip=10pt
}
\tcolorboxenvironment{obstruction}{
enhanced jigsaw,colframe=red,interior hidden,
breakable,%before skip=10pt,after skip=10pt
}

\tcolorboxenvironment{proof}{
enhanced jigsaw, frame hidden,%interior hidden,
breakable,%before skip=10pt,after skip=10pt
}

\lhead{Avi Chetlin \\ Assignment 01 \\ 31 August 2026}
\rhead{Dr. Stanislaw Szarek \\ MATH 423 Measure Theory \\ Fall 2026}

\begin{document}
For a set \( A \subset [a,b] \), define the outer Riemann measure
\begin{equation}
\label{eq:1}
m^{\star}(A) = \inf_{\cup J_i \subset A} \sum\limits_i l(J_i),
\end{equation}
where \( \left\{ J_i \right\} \) is a finite family of closed intervals contained in \( [a,b] \) and \( l(J) \) is the length of the interval \( J \).
One similarly defines the inner measure as
\begin{equation}
\label{eq:2}
m_{\star}(A) = \sup_{\cup J_i \subset A \\ J_i \text{disjoint}} \sum\limits_i l(J_i).
\end{equation}
If \( m^{\star}(A) = m_{\star}(A) \), then \( A \) is Riemann measureable and write \( A \in \mathcal{R}([a,b]) \).

\begin{problem}{1}
  Show that the value of \( m^{\star} \) remains unchanged if
  \begin{enumerate}
    \item we require that the intervals \( (J_i) \) in the definition \cref{eq:1} are disjoint
    \item we require that the family \( (J_i) \) is minimal (i.e., no proper subfamily of \( (J_i) \) covers \( A \)).
    \item we require that the intervals \( J_i \) are open (rather than closed).
  \end{enumerate}
\end{problem}

\begin{problem}{2}
  Show that \( m^{\star}(A) = U(\chi_A) \) (see the handout) where \(\chi_A\) is the indicator function of \( A \).
  \textit{Note}: This is relatively subtle;
  start by associating with a partition of \( [a,b] \) the family of intervals of that partition that intersect \( A \).
\end{problem}

\begin{problem}{3}
  One can show that \( m_{\star}(A) = L(\chi_A) \).
  Explain why this implies that \( A \) is Riemann measureable if and only if \(\chi_A \in \mathcal{R}[a,b]\), and that then \( m(A) = \int\limits_a^b \chi_A \).
\end{problem}

\begin{problem}{4}
  Show that if \( m^{\star}(A) = 0 \), then \( A \) is Riemann measureable.
\end{problem}

\begin{problem}{5}
  Suppose \(\mu: \mathcal{P}(\mathbb{R}) \rightarrow [0, \infty]\) is a finitely additive measure verifying
  \begin{enumerate}
    \item \(\mu(J) = l(J)\) for every interval \( J \subset \mathbb{R} \) and which is
    \item translation invariant.
  \end{enumerate}
  What is \(\mu(N)\), where \( N \subset [0,1) \) is the non-measureable set defined on p. 20 of the textbook?
  Recall that it was shown that \(\mu: \mathcal{P}(\mathbb{R}) \rightarrow [0, \infty]\) verifying (1) and (2) can not be countably additive, but the argument did not exclude the existence of such finitely additive measure.
\end{problem}

\end{document}
